\par\medskip\noindent\textbf{Cómo leer la salida.}
{\small\setlength{\tabcolsep}{4pt}\renewcommand{\arraystretch}{1.15}
\begin{xltabular}{\linewidth}{@{}>{\raggedright\arraybackslash}p{0.24\linewidth}>{\raggedleft\arraybackslash}p{0.15\linewidth}>{\raggedright\arraybackslash}X@{}}
\toprule Indicador & Valor & Interpretación \\ \midrule \endfirsthead
\toprule Indicador & Valor & Interpretación \\ \midrule \endhead
\bottomrule \endlastfoot
$n$ & 2, 3, 6, 12 & Vencimiento del bono que se compra y vende al mes siguiente. \\
\texttt{media rx} & 0,030 a 0,069 & Exceso de rendimiento medio mensual (\%) de mantener el bono un mes frente a la letra de un mes: crece con el vencimiento (prima por plazo). \\
\texttt{desv.est.} & 0,051 a 0,488 & Volatilidad del exceso: crece mucho más rápido que la media. \\
\texttt{Sharpe anual} & 2,05 a 0,49 & La razón de Sharpe cae con el vencimiento: los bonos cortos pagaban mucho por unidad de riesgo. \\
\texttt{b (Fama-Bliss)} & 0,52 a 1,30 & Coeficiente del diferencial a plazo; la HE predice 0. Que sea cercano a 1 indica que el diferencial predice el exceso de rendimiento, no la tasa futura. \\
\texttt{ee}, \texttt{t(b=0)} & $t=2{,}41$ a 5,21 & Significativos en los cuatro vencimientos: la prima por plazo varía en el tiempo. \\
\texttt{R2} & 0,035 a 0,118 & El diferencial explica entre el 3 y el 12\,\% del exceso mensual. \\
\end{xltabular}}
